\begin{table}[htbp]
\centering
\small
\caption{Problema G (google/gemini-2.5-flash-lite): nomes de classe semânticos versus neutros.}
\label{tab:bloco3_G_nomes_semanticos_vs_neutros_flashlite}
\begin{tabular}{llrrrrrr}
\toprule
 &  & \multicolumn{3}{c}{LLM vs real (\%)} & \multicolumn{3}{c}{Métrica vs real (\%)} \\ \cmidrule(lr){3-5} \cmidrule(lr){6-8}
Subproblema & Atributos & Mulher/Homem & A/B & $\Delta$ & Mulher/Homem & A/B & $\Delta$ \\
\midrule
G1 (normalizado) & peso/altura & 35,2 & 18,6 & +16,7 & 10,5 & 16,7 & $-$6,2 \\
G1 (normalizado) & x1/x2 & 57,1 & 16,2 & +41,0 & 74,8 & 9,5 & +65,2 \\
\midrule
G2 (kg e m) & peso\_kg/altura\_m & 84,3 & 46,7 & +37,6 & 91,4 & 22,9 & +68,6 \\
G2 (kg e m) & x1/x2 & 48,6 & — & — & 65,2 & — & — \\
\bottomrule
\multicolumn{8}{p{0.95\linewidth}}{\footnotesize Efeito do nome das classes (Mulher/Homem versus A/B) para cada par de nomes de atributo, zero-shot, R$^2$, Perceptron. $\Delta$ = semântico $-$ neutro. ``—'' = colapso em classe única.} \\
\end{tabular}
\end{table}
